%============================================================
%  Bd. 15 - Reelle Zahlen
%============================================================

\documentclass[main.tex]{subfiles}

\ifSubfilesClassLoaded{
  \BandSetupStandalone{B15}
}{
}

\title{Bd. 15 - Reelle Zahlen}
\author{Martin Kunze}
\date{}
\setcounter{file}{15}

\hypersetup{
  pdftitle={Bd. 15 - Reelle Zahlen},
  pdfauthor={Martin Kunze}
}

\begin{document}

\title{Bd. 15 - Reelle Zahlen}
\author{Martin Kunze}
\date{}
\hypersetup{pageanchor=false}
\maketitle
\hypersetup{pageanchor=true}
\setcounter{tocdepth}{3}
\tableofcontents
\setcounter{file}{15}
\renewcommand{\FormulaBandID}{15}

% Lokale Notation dieses Bandes. Die Definitionen stehen absichtlich innerhalb
% des Dokumentteils, damit sie sowohl im Gesamtband als auch im Standalone-Build
% zur Verfuegung stehen.
\providecommand{\RealCut}[1]{\operatorname{Cut}_{\mathbb Q}\!\left(#1\right)}
\providecommand{\RealHat}[1]{\widehat{#1}}
\providecommand{\RealEmb}{\iota_{\mathbb R}}
\providecommand{\RealLe}{\mathrel{\leq_{\mathbb R}}}
\providecommand{\RealLt}{\mathrel{<_{\mathbb R}}}
\providecommand{\RealZero}{0_{\mathbb R}}
\providecommand{\RealOne}{1_{\mathbb R}}
\providecommand{\RealAdd}{\mathbin{\oplus}}
\providecommand{\RealMul}{\mathbin{\odot}}
\providecommand{\RealPosMul}{\mathbin{\odot_{+}}}
\providecommand{\RealNeg}[1]{\mathop{\ominus}\!#1}
\providecommand{\RealPosInv}[1]{\operatorname{inv}_{+}\!\left(#1\right)}
\providecommand{\RealInv}[1]{\operatorname{inv}_{\mathbb R}\!\left(#1\right)}
\providecommand{\RealAbs}[1]{\left\lvert #1\right\rvert_{\mathbb R}}
\providecommand{\RealDist}[2]{d_{\mathbb R}\!\left(#1,#2\right)}
\providecommand{\CompleteOrderedField}[1]{\mathsf{COF}\!\left(#1\right)}

\chapter{Einleitung}

Band~14 konstruiert die rationalen Zahlen mit ihrer Addition,
Multiplikation und totalen Ordnung; die maßgeblichen Grundlagen sind in
\FormulaRefAuto{RationalOrderedFieldAxioms} zusammengefasst. In diesem Band
werden daraus die reellen Zahlen als Dedekindsche Schnitte gewonnen. Diese
Konstruktion verwendet nur Mengen, rationale Arithmetik und Ordnung;
insbesondere setzt sie weder Folgen noch einen Grenzwertbegriff voraus.

Ein reeller Wert wird durch die Menge aller rationalen Zahlen beschrieben,
die echt unter ihm liegen. Die Inklusion solcher Mengen liefert die Ordnung.
Das Supremum einer nichtleeren, nach oben beschränkten Familie ist ihre
Vereinigung. Damit ist die Vollständigkeit bereits auf der Ebene der
Konstruktion sichtbar. Folgen und ihre Grenzwerte können deshalb erst in
einem späteren Band ohne begrifflichen Zirkel eingeführt werden.

\chapter{Dedekindsche Schnitte}

\section{Schnittprädikat und Trägermenge}

\FormulaDefDeltaK[Dedekindscher Schnitt der rationalen Zahlen]{%
  \begin{aligned}[t]
  \RealCut{A}\coloneqq{}&
    A\subseteq\mathbb Q\land A\neq\varnothing\land A\neq\mathbb Q
    \land{}\\[-2pt]
  &\forall p\in A\,\forall q\in\mathbb Q\,
      \bigl(q<p\rightarrow q\in A\bigr)
    \land{}\\[-2pt]
  &\forall p\in A\,\exists r\in A\,(p<r).
  \end{aligned}%
}{DedekindCutPredicate}{%
  \DeltaRow{Mengen}{A\dsep p\dsep q\dsep r}%
  \DeltaRow{Relationsoperatoren}{<}%
  \DeltaPrem{Geordneter rationaler Körper}{\mathbb Q}%
  \DeltaRow{\textbf{Bedingungen}}{}%
  \DeltaRow{Nichtleer und echt}{A\neq\varnothing\dsep A\neq\mathbb Q}%
  \DeltaRow{Nach unten abgeschlossen}
    {p\in A\dsep q\in\mathbb Q\dsep q<p\vdash q\in A}%
  \DeltaRow{Kein größtes Element}
    {p\in A\vdash\exists r\in A\,(p<r)}%
  \DeltaRow{\textbf{Neues Symbol}}{}%
  \DeltaRow{Schnittprädikat}{\RealCut{\,\cdot\,}}%
}

Die vier Bedingungen sind unabhängig voneinander wichtig. Die
Abwärtsabgeschlossenheit macht einen Schnitt zu einem rationalen
Anfangsabschnitt; die letzte Bedingung entfernt die ansonsten mögliche
Mehrdeutigkeit am Rand eines rationalen Anfangsabschnitts.

\FormulaDefDeltaK[Die Menge der reellen Zahlen]{%
  \mathbb R\coloneqq
  \bigl\{A\in\powerset(\mathbb Q)\bigm|\RealCut{A}\bigr\}%
}{RealCutCarrier}{%
  \DeltaRow{Mengen}{A\dsep\mathbb R}%
  \DeltaRow{Prädikate}{\RealCut{\,\cdot\,}}%
  \DeltaRow{\textbf{Neues Symbol}}{}%
  \DeltaRow{Reelle Zahlen}{\mathbb R}%
}

\FormulaThmDeltaK[Elementcharakterisierung der reellen Zahlen]{%
  A\in\mathbb R\eqvdash\RealCut{A}%
}{RealCutMembership}{%
  \DeltaRow{Mengen}{A}%
}

Diese Charakterisierung erlaubt es im Folgenden, eine reelle Zahl und den sie
darstellenden Schnitt ohne zusätzliche Klassenklammern miteinander zu
identifizieren.

\section{Hauptschnitte und rationale Einbettung}

\FormulaDefDeltaK[Hauptschnitt einer rationalen Zahl]{%
  q\in\mathbb Q\vdash
  \RealHat q\coloneqq\{p\in\mathbb Q\mid p<q\}%
}{RealPrincipalCut}{%
  \DeltaRow{Mengen}{p\dsep q}%
  \DeltaRow{Relationsoperatoren}{<}%
  \DeltaRow{\textbf{Neues Symbol}}{}%
  \DeltaRow{Hauptschnitt}{\RealHat q}%
}

\FormulaThmDeltaK[Hauptschnitte sind Dedekindsche Schnitte]{%
  q\in\mathbb Q\vdash\RealCut{\RealHat q}%
}{RealPrincipalCutIsCut}{%
  \DeltaRow{Mengen}{p\dsep q\dsep r}%
}

Der Nachweis verwendet die Unbeschränktheit und die in
\FormulaRefAuto{RationalOrderDense}[thm] registrierte Dichtheit der rationalen
Ordnung aus Band~14: Unterhalb von \(q\) liegt eine rationale Zahl, oberhalb
jedes \(p<q\) liegt noch eine rationale Zahl \(r<q\), und \(q\) selbst gehört
nicht zu \(\RealHat q\).

\FormulaDefDeltaK[Kanonische Einbettung der rationalen Zahlen]{%
  q\in\mathbb Q\vdash\RealEmb(q)\coloneqq\RealHat q%
}{RationalToRealEmbedding}{%
  \DeltaRow{Mengen}{p\dsep q}%
  \DeltaRow{Funktionen}{\RealEmb}%
  \DeltaRow{\textbf{Neues Symbol}}{}%
  \DeltaRow{Rationale Einbettung}{\RealEmb}%
}

\FormulaThmDeltaK[Die rationale Einbettung ist eine Funktion]{%
  \RealEmb\colon\mathbb Q\to\mathbb R%
}{RationalToRealEmbeddingFunction}{%
  \DeltaRow{Funktionen}{\RealEmb}%
}

\chapter{Ordnung und Vollständigkeit}

\section{Inklusionsordnung}

\FormulaDefDeltaK[Ordnung der reellen Zahlen]{%
  A,B\in\mathbb R\vdash
  A\RealLe B\coloneqq A\subseteq B%
}{RealCutOrder}{%
  \DeltaRow{Mengen}{A\dsep B}%
  \DeltaRow{Relationsoperatoren}{\RealLe}%
  \DeltaRow{\textbf{Neues Symbol}}{}%
  \DeltaRow{Reelle Ordnung}{\RealLe}%
}

\FormulaDefDeltaK[Strikte Ordnung der reellen Zahlen]{%
  A,B\in\mathbb R\vdash
  A\RealLt B\coloneqq A\subsetneq B%
}{RealCutStrictOrder}{%
  \DeltaRow{Mengen}{A\dsep B}%
  \DeltaRow{Relationsoperatoren}{\RealLe\dsep\RealLt}%
  \DeltaRow{\textbf{Neues Symbol}}{}%
  \DeltaRow{Strikte reelle Ordnung}{\RealLt}%
}

\FormulaThmDeltaK[Charakterisierung der reellen Vergleichsrelationen]{%
  A,B\in\mathbb R\vdash
  \bigl(A\RealLe B\leftrightarrow A\subseteq B\bigr)
  \land
  \bigl(A\RealLt B\leftrightarrow A\subsetneq B\bigr)%
}{RealCutOrderCharacterization}{%
  \DeltaRow{Mengen}{A\dsep B}%
  \DeltaRow{Relationsoperatoren}{\RealLe\dsep\RealLt}%
}

\FormulaThmDeltaK[Die Inklusionsordnung der Schnitte ist total]{%
  \TotOrd{\mathbb R,\RealLe}%
}{RealCutTotalOrder}{%
  \DeltaRow{Mengen}{A\dsep B}%
  \DeltaRow{Relationsoperatoren}{\RealLe}%
}

Für die Vergleichbarkeit seien \(a\in A\setminus B\) und
\(b\in B\setminus A\) angenommen. Aus \(a<b\) folgte wegen der
Abwärtsabgeschlossenheit \(a\in B\), aus \(b<a\) entsprechend \(b\in A\);
auch \(a=b\) ist unmöglich. Die Totalität der rationalen Ordnung schließt
daher aus, dass beide Mengendifferenzen zugleich nichtleer sind.

\FormulaThmDeltaK[Hauptschnitte bewahren und spiegeln die Ordnung]{%
  p,q\in\mathbb Q\vdash
  \bigl(p\leq q\leftrightarrow\RealEmb(p)\RealLe\RealEmb(q)\bigr)
  \land
  \bigl(p<q\leftrightarrow\RealEmb(p)\RealLt\RealEmb(q)\bigr)%
}{RationalToRealEmbeddingOrder}{%
  \DeltaRow{Mengen}{p\dsep q}%
  \DeltaRow{Funktionen}{\RealEmb}%
  \DeltaRow{Relationsoperatoren}{\RealLe\dsep\RealLt}%
}

\FormulaThmDeltaK[Injektivität der rationalen Einbettung]{%
  \RealEmb\colon\mathbb Q\inj\mathbb R%
}{RationalToRealEmbeddingInjective}{%
  \DeltaRow{Funktionen}{\RealEmb}%
}

\section{Suprema als Vereinigungen}

\FormulaThmDeltaK[Vereinigung einer beschränkten Schnittfamilie]{%
  \begin{aligned}[t]
    &\varnothing\neq\mathcal A\subseteq\mathbb R\dsep
      U\in\mathbb R\dsep
      \forall A\in\mathcal A\,(A\RealLe U)\\[-2pt]
    &\hspace{34mm}\vdash\RealCut{\bigcup\mathcal A}.
  \end{aligned}%
}{RealBoundedCutUnionIsCut}{%
  \DeltaRow{Mengen}{\mathcal A\dsep A\dsep U}%
  \DeltaRow{Relationsoperatoren}{\RealLe}%
}

Die Nichtleerheit eines Mitglieds von \(\mathcal A\) liefert die
Nichtleerheit der Vereinigung. Aus
\(\bigcup\mathcal A\subseteq U\neq\mathbb Q\) folgt ihre Echtheit. Die beiden
Abschlussbedingungen werden jeweils in dem Schnitt nachgewiesen, aus dem das
betrachtete Element der Vereinigung stammt.

\FormulaThmDeltaK[Vollständigkeit der reellen Ordnung]{%
  \begin{aligned}[t]
    &\varnothing\neq\mathcal A\subseteq\mathbb R\dsep
      \exists U\in\mathbb R\,
        \forall A\in\mathcal A\,(A\RealLe U)\\[-2pt]
    &\hspace{19mm}\vdash
      \exists S\in\mathbb R\,
      \left(
        \begin{aligned}[t]
        &\forall A\in\mathcal A\,(A\RealLe S)\land{}\\[-2pt]
        &\forall V\in\mathbb R\,
          \bigl(\forall A\in\mathcal A\,(A\RealLe V)
          \rightarrow S\RealLe V\bigr)
        \end{aligned}
      \right).
  \end{aligned}%
}{RealLeastUpperBoundProperty}{%
  \DeltaRow{Mengen}{\mathcal A\dsep A\dsep S\dsep U\dsep V}%
  \DeltaRow{Relationsoperatoren}{\RealLe}%
}

\emph{Beweisskizze.}
Wähle eine obere Schranke \(U\) und setze \(S=\bigcup\mathcal A\). Nach
\FormulaRefAuto{RealBoundedCutUnionIsCut}[thm] gilt \(S\in\mathbb R\).
Jedes \(A\in\mathcal A\) ist in \(S\) enthalten. Ist \(V\in\mathbb R\)
eine weitere obere Schranke, so liegt jedes Element der Vereinigung in einem
\(A\in\mathcal A\) und damit in \(V\). Also ist \(S\subseteq V\), und \(S\)
ist die kleinste obere Schranke.

Erst die vorstehende Existenz- und Kleinste-Obergrenzen-Aussage erfüllt die
Definitionsprämisse des allgemeinen Supremumsbegriffs aus Band~11. Daher ist
der folgende Supremumsterm wohldefiniert.

\FormulaThmDeltaK[Supremum einer beschränkten Schnittfamilie]{%
  \begin{aligned}[t]
    &\varnothing\neq\mathcal A\subseteq\mathbb R\dsep
      U\in\mathbb R\dsep
      \forall A\in\mathcal A\,(A\RealLe U)\\[-2pt]
    &\hspace{20mm}\vdash
      \sup_{\mathbb R,\RealLe}(\mathcal A)=\bigcup\mathcal A.
  \end{aligned}%
}{RealCutSupremumUnion}{%
  \DeltaRow{Mengen}{\mathcal A\dsep A\dsep U}%
  \DeltaRow{Relationsoperatoren}{\RealLe}%
}

\emph{Beweisskizze.}
Die vorherige Aussage liefert die Existenz einer kleinsten oberen Schranke;
damit ist das Supremum nach Band~11 definiert. Der dortige
Charakterisierungssatz und der eben geführte Vereinigungsbeweis identifizieren
diese eindeutig bestimmte Schranke mit \(\bigcup\mathcal A\).

\FormulaThmDeltaK[Infimumsvollständigkeit der reellen Ordnung]{%
  \begin{aligned}[t]
    &\varnothing\neq\mathcal A\subseteq\mathbb R\dsep
      \exists L\in\mathbb R\,
        \forall A\in\mathcal A\,(L\RealLe A)\\[-2pt]
    &\quad\vdash \exists I\in\mathbb R\,
      \left(
        \forall A\in\mathcal A\,(I\RealLe A)\land
        \forall V\in\mathbb R\,
          \bigl(\forall A\in\mathcal A\,(V\RealLe A)
          \rightarrow V\RealLe I\bigr)
      \right).
  \end{aligned}%
}{RealGreatestLowerBoundProperty}{%
  \DeltaRow{Mengen}{\mathcal A\dsep A\dsep L\dsep V}%
  \DeltaRow{Relationsoperatoren}{\RealLe}%
}

\emph{Beweisskizze.}
Die Menge aller unteren Schranken von \(\mathcal A\) ist wegen \(L\)
nichtleer und wird durch jedes feste \(A_0\in\mathcal A\) nach oben
beschränkt. Ihr Supremum existiert nach
\FormulaRefAuto{RealLeastUpperBoundProperty}[thm]. Jede untere Schranke liegt
darunter. Da jedes \(A\in\mathcal A\) eine obere Schranke der Menge aller
unteren Schranken ist, liegt deren Supremum seinerseits unter jedem \(A\) und
ist somit die größte untere Schranke.

Auch hier wird der allgemeine Infimumsterm erst nach diesem Existenzsatz
eingeführt.

\FormulaThmDeltaK[Charakterisierung des reellen Infimums]{%
  \begin{aligned}[t]
    &\varnothing\neq\mathcal A\subseteq\mathbb R\dsep
      \exists L\in\mathbb R\,
        \forall A\in\mathcal A\,(L\RealLe A)\\[-2pt]
    &\quad\vdash
      \inf_{\mathbb R,\RealLe}(\mathcal A)\in\mathbb R\land{}\\[-2pt]
    &\qquad
      \forall A\in\mathcal A\,
        \bigl(\inf_{\mathbb R,\RealLe}(\mathcal A)\RealLe A\bigr)\land{}\\[-2pt]
    &\qquad
      \forall V\in\mathbb R\,
        \bigl(\forall A\in\mathcal A\,(V\RealLe A)
        \rightarrow
        V\RealLe\inf_{\mathbb R,\RealLe}(\mathcal A)\bigr).
  \end{aligned}%
}{RealInfimumCharacterization}{%
  \DeltaRow{Mengen}{\mathcal A\dsep A\dsep L\dsep V}%
  \DeltaRow{Relationsoperatoren}{\RealLe}%
}

\emph{Beweisskizze.}
\FormulaRefAuto{RealGreatestLowerBoundProperty}[thm] liefert einen Kandidaten,
der in der Menge aller unteren Schranken maximal ist. Damit ist die
Definitionsprämisse von \FormulaRefAuto{Infimum}[def] erfüllt; deren
Charakterisierung liefert die drei angezeigten Aussagen.

\chapter{Arithmetik der Schnitte}

\section{Addition und additive Inverse}

\FormulaDefDeltaK[Addition reeller Schnitte]{%
  A,B\in\mathbb R\vdash
  A\RealAdd B\coloneqq
  \bigl\{q\in\mathbb Q\bigm|
    \exists a\in A\,\exists b\in B\,(q<a+b)\bigr\}%
}{RealCutAddition}{%
  \DeltaRow{Mengen}{A\dsep B\dsep a\dsep b\dsep q}%
  \DeltaRow{Funktionsoperatoren}{\RealAdd}%
  \DeltaRow{\textbf{Neues Symbol}}{}%
  \DeltaRow{Reelle Addition}{\RealAdd}%
}

\FormulaDefDeltaK[Additive Null der reellen Zahlen]{%
  \RealZero\coloneqq\RealHat{0_{\mathbb Q}}%
}{RealCutZero}{%
  \DeltaRow{Mengen}{\RealZero}%
  \DeltaRow{\textbf{Neues Symbol}}{}%
  \DeltaRow{Reelle Null}{\RealZero}%
}

\FormulaDefDeltaK[Additives Inverses eines reellen Schnittes]{%
  A\in\mathbb R\vdash
  \RealNeg{A}\coloneqq
  \bigl\{q\in\mathbb Q\bigm|
    \exists r\in\mathbb Q\,(r\notin A\land q<-r)\bigr\}%
}{RealCutNegation}{%
  \DeltaRow{Mengen}{A\dsep q\dsep r}%
  \DeltaRow{Funktionsoperatoren}{\RealNeg{\,\cdot\,}}%
  \DeltaRow{\textbf{Neues Symbol}}{}%
  \DeltaRow{Reelle Negation}{\RealNeg{\,\cdot\,}}%
}

\FormulaThmDeltaK[Abgeschlossenheit der additiven Schnittoperationen]{%
  A,B\in\mathbb R\vdash
  A\RealAdd B\in\mathbb R\land
  \RealNeg{A}\in\mathbb R\land
  \RealZero\in\mathbb R%
}{RealCutAdditiveClosure}{%
  \DeltaRow{Mengen}{A\dsep B}%
  \DeltaRow{Funktionsoperatoren}{\RealAdd\dsep\RealNeg{\,\cdot\,}}%
}

\emph{Beweisskizze.}
Wähle \(a_0\in A,b_0\in B\) sowie
\(\alpha\in\mathbb Q\setminus A,\beta\in\mathbb Q\setminus B\). Eine
rationale Zahl unter \(a_0+b_0\) zeigt die Nichtleerheit von
\(A\RealAdd B\). Da jedes \(a\in A,b\in B\) die Ungleichungen
\(a\leq\alpha,b\leq\beta\) erfüllt, gehört \(\alpha+\beta\) nicht zur Summe;
sie ist also echt. Abwärtsabschluss ist unmittelbar, und zwischen einem
Element \(q\) der Summe und seinem bezeugenden Wert \(a+b\) liegt nach
\FormulaRefAuto{RationalOrderDense}[thm] ein größeres Summenelement.

Für \(\RealNeg A\) liefert ein \(r\notin A\) zusammen mit einer rationalen
Zahl unter \(-r\) die Nichtleerheit. Ist \(a\in A\), so kann \(-a\) nicht in
\(\RealNeg A\) liegen: Aus \(-a<-r\) folgte \(r<a\) und damit \(r\in A\).
Abwärtsabschluss und das Fehlen eines größten Elements folgen wieder direkt
aus der strikten Zeugenungleichung und rationaler Dichtheit. Damit erfüllen
beide Mengen \FormulaRefAuto{DedekindCutPredicate}[def]; für \(\RealZero\)
gilt dies nach \FormulaRefAuto{RealPrincipalCutIsCut}[thm].

\FormulaThmDeltaK[Rationale Randapproximation eines Schnittes]{%
  \begin{aligned}[t]
  &A\in\mathbb R\dsep h\in\mathbb Q\dsep0_{\mathbb Q}<h\vdash{}\\[-2pt]
  &\quad\exists a\in A\,\exists r\in\mathbb Q\,
    \bigl(r\notin A\land a<r\land r-a<h\bigr).
  \end{aligned}%
}{RealCutAdditiveBoundaryApproximation}{%
  \DeltaRow{Mengen}{A\dsep a\dsep r\dsep h}%
}

\emph{Beweisskizze.}
Wähle \(p\in A\) und \(u\in\mathbb Q\setminus A\). Das archimedische
Gesetz von \(\mathbb Q\) liefert für die positive Maschenweite \(h/2\) einen
Index \(m\in\mathbb N\), für den \(p+m(h/2)>u\) gilt. Unter allen solchen
Indizes, deren Gitterpunkt nicht in \(A\) liegt, sei \(m_0\) der kleinste.
Dann ist \(m_0>0\), der Vorgängerpunkt
\(a=p+(m_0-1)(h/2)\) liegt in \(A\), und
\(r=p+m_0(h/2)\) liegt nicht in \(A\). Somit
\(0<r-a=h/2<h\).

\FormulaThmDeltaK[Gesetze der reellen Addition]{%
  \begin{aligned}[t]
  A,B,C\in\mathbb R\vdash{}&
    (A\RealAdd B)\RealAdd C=A\RealAdd(B\RealAdd C)\dsep{}\\[-2pt]
  &A\RealAdd B=B\RealAdd A\dsep
    A\RealAdd\RealZero=A\dsep
    A\RealAdd\RealNeg{A}=\RealZero.
  \end{aligned}%
}{RealCutAdditiveLaws}{%
  \DeltaRow{Mengen}{A\dsep B\dsep C}%
  \DeltaRow{Funktionsoperatoren}{\RealAdd\dsep\RealNeg{\,\cdot\,}}%
}

\emph{Beweisskizze.}
Nach Auflösen von \FormulaRefAuto{RealCutAddition}[def] besitzen beide Seiten
der Assoziativität genau die rationalen Elemente, die echt unter einem Wert
\(a+b+c\) mit \(a\in A,b\in B,c\in C\) liegen; Kommutativität folgt ebenso
aus der rationalen Addition. Für die Null gilt
\(q\in A\RealAdd\RealZero\) genau dann, wenn \(q<a+r\) für geeignete
\(a\in A\) und \(r<0_{\mathbb Q}\); Abwärtsabschluss und die Existenz größerer
Elemente in \(A\) liefern beide Inklusionen. Schließlich trennt der Rand von
\(A\) die Summanden in den Definitionen von \(A\) und \(\RealNeg A\); daraus
folgt \(A\RealAdd\RealNeg A\subseteq\RealZero\). Für \(q<0_{\mathbb Q}\)
wähle rational \(0_{\mathbb Q}<h<-q\). Nach
\FormulaRefAuto{RealCutAdditiveBoundaryApproximation}[thm] gibt es
\(a\in A\) und \(r\notin A\) mit \(0<r-a<h\). Wegen \(q<a-r\) wählt man
\(b<-r\) noch so nahe an \(-r\), dass \(q<a+b\). Dann ist
\(b\in\RealNeg A\), also
\(q\in A\RealAdd\RealNeg A\), und die umgekehrte Inklusion ist gezeigt.

\FormulaThmDeltaK[Negation kehrt die Schnittordnung um]{%
  A,B\in\mathbb R\vdash
  \begin{aligned}[t]
    &\bigl(A\RealLe B\leftrightarrow
      \RealNeg B\RealLe\RealNeg A\bigr)\dsep
      \bigl(A\RealLt B\leftrightarrow
      \RealNeg B\RealLt\RealNeg A\bigr)\dsep{}\\[-2pt]
    &\RealNeg{\RealNeg A}=A\dsep
      \RealNeg{\RealZero}=\RealZero.
  \end{aligned}%
}{RealCutNegationOrder}{%
  \DeltaRow{Mengen}{A\dsep B}%
  \DeltaRow{Relationsoperatoren}{\RealLe\dsep\RealLt}%
  \DeltaRow{Funktionsoperatoren}{\RealNeg{\,\cdot\,}}%
}

\emph{Beweisskizze.}
Aus \(A\subseteq B\) folgt
\(\mathbb Q\setminus B\subseteq\mathbb Q\setminus A\); Einsetzen in
\FormulaRefAuto{RealCutNegation}[def] gibt die umgekehrte Inklusion der
Negationen. Zweimaliges Komplementieren der Randbedingung und die
No-Maximum-Bedingung liefern \(\RealNeg{\RealNeg A}=A\); für den Hauptschnitt
bei \(0_{\mathbb Q}\) ergibt dieselbe Definition
\(\RealNeg{\RealZero}=\RealZero\). Involutivität liefert die Rückrichtung der
nichtstrikten Ordnungsäquivalenz, und zusammen mit Ungleichheit folgt die
strikte Form.

\section{Multiplikation und multiplikative Inverse}

Zunächst wird das Produkt nichtnegativer Schnitte definiert. Die rationale
Null wird ausdrücklich beigefügt, damit auch das Produkt mit
\(\RealZero\) wieder der Schnitt \(\RealZero\) ist.

\FormulaDefDeltaK[Positives Produkt reeller Schnitte]{%
  \begin{aligned}[t]
  &A,B\in\mathbb R\dsep
    \RealZero\RealLe A\dsep\RealZero\RealLe B\\[-2pt]
  &\quad\vdash
  A\RealPosMul B\coloneqq
  \RealZero\cup
  \bigl\{q\in\mathbb Q\bigm|
    \exists a\in A\,\exists b\in B\,
      (0_{\mathbb Q}<a\land0_{\mathbb Q}<b\land q<ab)\bigr\}.
  \end{aligned}%
}{RealCutPositiveProduct}{%
  \DeltaRow{Mengen}{A\dsep B\dsep a\dsep b\dsep q}%
  \DeltaRow{Relationsoperatoren}{\RealLe}%
  \DeltaRow{Funktionsoperatoren}{\RealPosMul}%
  \DeltaRow{\textbf{Neues Symbol}}{}%
  \DeltaRow{Positives Schnittprodukt}{\RealPosMul}%
}

\FormulaDefDeltaK[Multiplikation beliebiger reeller Schnitte]{%
  A,B\in\mathbb R\vdash
  A\RealMul B\coloneqq
  \begin{cases}
    A\RealPosMul B,
      &\RealZero\RealLe A\land\RealZero\RealLe B,\\
    \RealNeg{\bigl((\RealNeg A)\RealPosMul B\bigr)},
      &A\RealLt\RealZero\land\RealZero\RealLe B,\\
    \RealNeg{\bigl(A\RealPosMul(\RealNeg B)\bigr)},
      &\RealZero\RealLe A\land B\RealLt\RealZero,\\
    (\RealNeg A)\RealPosMul(\RealNeg B),
      &A\RealLt\RealZero\land B\RealLt\RealZero.
  \end{cases}%
}{RealCutMultiplication}{%
  \DeltaRow{Mengen}{A\dsep B}%
  \DeltaRow{Relationsoperatoren}{\RealLe\dsep\RealLt}%
  \DeltaRow{Funktionsoperatoren}{\RealMul\dsep\RealPosMul
    \dsep\RealNeg{\,\cdot\,}}%
  \DeltaRow{\textbf{Neues Symbol}}{}%
  \DeltaRow{Reelle Multiplikation}{\RealMul}%
}

\FormulaDefDeltaK[Multiplikative Eins der reellen Zahlen]{%
  \RealOne\coloneqq\RealHat{1_{\mathbb Q}}%
}{RealCutOne}{%
  \DeltaRow{Mengen}{\RealOne}%
  \DeltaRow{\textbf{Neues Symbol}}{}%
  \DeltaRow{Reelle Eins}{\RealOne}%
}

\FormulaDefDeltaK[Positives Reziprokes eines Schnittes]{%
  \begin{aligned}[t]
  &A\in\mathbb R\dsep\RealZero\RealLt A\\[-2pt]
  &\quad\vdash
  \RealPosInv A\coloneqq
  \RealZero\cup
  \bigl\{q\in\mathbb Q\bigm|
    \exists r\in\mathbb Q\,
      (0_{\mathbb Q}<r\land r\notin A\land q<r^{-1})\bigr\}.
  \end{aligned}%
}{RealCutPositiveReciprocal}{%
  \DeltaRow{Mengen}{A\dsep q\dsep r}%
  \DeltaRow{Relationsoperatoren}{\RealLt}%
  \DeltaRow{Funktionsoperatoren}{\RealPosInv{\,\cdot\,}}%
  \DeltaRow{\textbf{Neues Symbol}}{}%
  \DeltaRow{Positives Reziprokes}{\RealPosInv{\,\cdot\,}}%
}

\FormulaDefDeltaK[Reziprokes einer von Null verschiedenen reellen Zahl]{%
  A\in\mathbb R\dsep A\neq\RealZero\vdash
  \RealInv A\coloneqq
  \begin{cases}
    \RealPosInv A,&\RealZero\RealLt A,\\
    \RealNeg{\bigl(\RealPosInv{\RealNeg A}\bigr)},&A\RealLt\RealZero.
  \end{cases}%
}{RealCutReciprocal}{%
  \DeltaRow{Mengen}{A}%
  \DeltaRow{Relationsoperatoren}{\RealLt}%
  \DeltaRow{Funktionsoperatoren}{\RealInv{\,\cdot\,}
    \dsep\RealPosInv{\,\cdot\,}\dsep\RealNeg{\,\cdot\,}}%
  \DeltaRow{\textbf{Neues Symbol}}{}%
  \DeltaRow{Reelles Reziprokes}{\RealInv{\,\cdot\,}}%
}

\FormulaThmDeltaK[Abgeschlossenheit der multiplikativen Schnittoperationen]{%
  \begin{aligned}[t]
    &A,B\in\mathbb R\dsep
      \RealZero\RealLe A\dsep\RealZero\RealLe B
      \vdash A\RealPosMul B\in\mathbb R,\\[-2pt]
    &A,B\in\mathbb R
      \vdash A\RealMul B\in\mathbb R\land\RealOne\in\mathbb R,\\[-2pt]
    &A\in\mathbb R\dsep\RealZero\RealLt A
      \vdash\RealPosInv A\in\mathbb R,\\[-2pt]
    &A\in\mathbb R\dsep A\neq\RealZero
      \vdash\RealInv A\in\mathbb R.
  \end{aligned}%
}{RealCutMultiplicativeClosure}{%
  \DeltaRow{Mengen}{A\dsep B}%
  \DeltaRow{Funktionsoperatoren}{\RealPosMul\dsep\RealMul
    \dsep\RealPosInv{\,\cdot\,}\dsep\RealInv{\,\cdot\,}}%
}

\emph{Beweisskizze.}
Sind \(A,B\) nichtnegativ, so enthält \(A\RealPosMul B\) den Schnitt
\(\RealZero\). Ist ein Faktor gleich null, ist der positive Teil leer und das
Produkt gleich \(\RealZero\). Andernfalls besitzen beide Faktoren positive
Elemente und positive rationale Nichtmitglieder \(\alpha\notin A\),
\(\beta\notin B\). Dann liegt \(\alpha\beta\) nicht im Produkt, weil
\(0_{\mathbb Q}<a\in A,0_{\mathbb Q}<b\in B\) stets
\(a\leq\alpha,b\leq\beta\) erfüllen. Dies zeigt
die Echtheit. Abwärtsabschluss und das Fehlen eines größten Elements folgen
aus der strikten Ungleichung \(q<ab\) und rationaler Dichtheit.

Nach \FormulaRefAuto{RealCutAdditiveClosure}[thm] sind die Negationen
Schnitte; \FormulaRefAuto{RealCutNegationOrder}[thm] macht einen negativen
Faktor nach Negation positiv. Totalität der Ordnung macht die vier
Vorzeichenfälle disjunkt und vollständig und typisiert somit jeden Zweig von
\FormulaRefAuto{RealCutMultiplication}[def].

Für \(\RealZero\RealLt A\) wähle \(0_{\mathbb Q}<a\in A\). Das positive
Komplement von \(A\) ist nichtleer, und jedes seiner Elemente \(r\) erfüllt
\(a\leq r\). Folglich gehört \(a^{-1}\) nicht zu \(\RealPosInv A\); der
Kandidat ist also echt, während \(\RealZero\) seine Nichtleerheit sichert.
Abwärtsabschluss und das Fehlen eines größten Elements folgen erneut aus der
strikten Zeugenungleichung \(q<r^{-1}\) und rationaler Dichtheit. Der negative
Fall folgt durch Negation. Die Eins ist nach
\FormulaRefAuto{RealPrincipalCutIsCut}[thm] ein Schnitt.

\FormulaThmDeltaK[Multiplikative Randapproximation eines positiven Schnittes]{%
  \begin{aligned}[t]
  &A\in\mathbb R\dsep\RealZero\RealLt A\dsep
    q\in\mathbb Q\dsep0_{\mathbb Q}<q<1_{\mathbb Q}\vdash{}\\[-2pt]
  &\quad\exists a\in A\,\exists r\in\mathbb Q\,
    \bigl(0_{\mathbb Q}<a<r\land r\notin A\land q<a r^{-1}\bigr).
  \end{aligned}%
}{RealCutMultiplicativeBoundaryApproximation}{%
  \DeltaRow{Mengen}{A\dsep a\dsep r\dsep q}%
  \DeltaRow{Relationsoperatoren}{\RealLt}%
}

\emph{Beweisskizze.}
Aus \(\RealZero\RealLt A\) und dem Fehlen eines größten Schnittpunkts folgt
die Existenz eines \(p\in A\) mit \(0_{\mathbb Q}<p\). Wähle rational
\(h>0\) so klein, dass \(qh<(1-q)p\). Wiederhole den archimedischen
Gitterbeweis mit dem fest vorgegebenen Startpunkt \(p\): Unter den Indizes
\(m\in\mathbb N\), für die \(p+mh\notin A\) gilt, existiert wegen eines
positiven rationalen Nichtmitglieds von \(A\) und der Archimedizität ein
kleinster Index \(m_0>0\). Mit
\(a=p+(m_0-1)h\) und \(r=p+m_0h\) erhält man daher
\(p\leq a<r=a+h\), \(a\in A\) und \(r\notin A\). Dann
\[
  qr=qa+qh<qa+(1-q)p\leq qa+(1-q)a=a,
\]
also \(q<a/r\), wie behauptet.

\FormulaThmDeltaK[Gesetze der reellen Multiplikation]{%
  \begin{aligned}[t]
  A,B,C\in\mathbb R\vdash{}&
    (A\RealMul B)\RealMul C=A\RealMul(B\RealMul C)\dsep{}\\[-2pt]
  &A\RealMul B=B\RealMul A\dsep
    A\RealMul\RealOne=A\dsep{}\\[-2pt]
  &A\RealMul(B\RealAdd C)
      =(A\RealMul B)\RealAdd(A\RealMul C),\\[-2pt]
  &A\neq\RealZero\rightarrow A\RealMul\RealInv A=\RealOne.
  \end{aligned}%
}{RealCutMultiplicativeLaws}{%
  \DeltaRow{Mengen}{A\dsep B\dsep C}%
  \DeltaRow{Funktionsoperatoren}{\RealAdd\dsep\RealMul
    \dsep\RealInv{\,\cdot\,}}%
}

\emph{Beweisskizze.}
Für nichtnegative Schnitte werden Gleichheiten durch beidseitige Inklusion
gezeigt. Nach Auflösen der Definitionen reduzieren Assoziativität,
Kommutativität, Einheitsgesetz und Distributivität auf die entsprechenden
rationalen Gesetze; die jeweils strikten Randungleichungen werden mit der
Dichtheit von \(\mathbb Q\) angepasst. Für das positive Reziproke zeigt
\FormulaRefAuto{RealCutMultiplicativeBoundaryApproximation}[thm] zu jedem
\(0<q<1_{\mathbb Q}\) genau die Zeugen \(a\in A\), \(r\notin A\) mit
\(q/a<r^{-1}\). Nach rationaler Dichtheit wähle
\(b\) mit \(q/a<b<r^{-1}\). Dann ist
\(b\in\RealPosInv A\) und \(q<ab\), also
\(q\in A\RealPosMul\RealPosInv A\). Umgekehrt besitzt jedes positive
\(x\in A\RealPosMul\RealPosInv A\) Zeugen \(a\in A\),
\(b<r^{-1}\) und \(r\notin A\) mit \(x<ab\). Wegen \(a\leq r\) folgt
\(x<ab<a/r\leq1_{\mathbb Q}\). Somit gilt
\(A\RealPosMul\RealPosInv A=\RealOne\).
Die Vorzeichenzerlegung aus
\FormulaRefAuto{RealCutMultiplication}[def] gibt unmittelbar die
Vorzeichenidentitäten
\((\RealNeg A)\RealMul B=\RealNeg{(A\RealMul B)}\) und
\((\RealNeg A)\RealMul(\RealNeg B)=A\RealMul B\). Eine endliche
Fallunterscheidung überträgt Assoziativität und die Inverseigenschaft auf
beliebige Schnitte. Für die Distributivität bleibt nur der Fall verschieden
gerichteter Vorzeichen von \(B\) und \(C\): Ist etwa
\(B\geq0>C\) und \(B+C\geq0\), so gilt
\(B=(B+C)+(-C)\); positive Distributivität und additives Umstellen liefern
die gewünschte Gleichung. Ist \(B+C<0\), verwendet man stattdessen
\(-C=B+(-(B+C))\). Der Fall eines negativen Faktors \(A\) folgt aus der
ersten Vorzeichenidentität. Damit sind auch die gemischten Fälle erfasst.

\FormulaThmDeltaK[Verträglichkeit von Ordnung und Arithmetik]{%
  \begin{aligned}[t]
  A,B,C\in\mathbb R\vdash{}&
    \bigl(A\RealLe B\rightarrow
      A\RealAdd C\RealLe B\RealAdd C\bigr)\land{}\\[-2pt]
  &\bigl(\RealZero\RealLe A\land\RealZero\RealLe B
      \rightarrow\RealZero\RealLe A\RealMul B\bigr).
  \end{aligned}%
}{RealCutOrderArithmeticCompatibility}{%
  \DeltaRow{Mengen}{A\dsep B\dsep C}%
  \DeltaRow{Relationsoperatoren}{\RealLe}%
  \DeltaRow{Funktionsoperatoren}{\RealAdd\dsep\RealMul}%
}

\emph{Beweisskizze.}
Aus \(A\subseteq B\) bleibt jeder Zeuge \(a\in A,c\in C\) auch ein Zeuge für
\(B\RealAdd C\); dies gibt die Translationsmonotonie als Mengeninklusion. Sind
beide Faktoren nichtnegativ, verwendet ihre Multiplikation unmittelbar
\FormulaRefAuto{RealCutPositiveProduct}[def] und enthält daher
\(\RealZero\).

\section{Erhaltung der rationalen Arithmetik}

\FormulaThmDeltaK[Die rationale Einbettung erhält die Körperoperationen]{%
  \begin{aligned}[t]
  p,q\in\mathbb Q\vdash{}&
    \RealEmb(p+q)=\RealEmb(p)\RealAdd\RealEmb(q)\dsep{}\\[-2pt]
  &\RealEmb(pq)=\RealEmb(p)\RealMul\RealEmb(q)\dsep{}\\[-2pt]
  &\RealEmb(-p)=\RealNeg{\RealEmb(p)}\dsep{}\\[-2pt]
  &p\neq0_{\mathbb Q}\rightarrow
    \RealEmb(p^{-1})=\RealInv{\RealEmb(p)}.
  \end{aligned}%
}{RationalToRealEmbeddingArithmetic}{%
  \DeltaRow{Mengen}{p\dsep q}%
  \DeltaRow{Funktionen}{\RealEmb}%
  \DeltaRow{Funktionsoperatoren}{\RealAdd\dsep\RealMul
    \dsep\RealNeg{\,\cdot\,}\dsep\RealInv{\,\cdot\,}}%
}

\emph{Beweisskizze.}
Für Hauptschnitte werden die vier Gleichheiten durch Einsetzen der
Definitionen bewiesen. Beispielsweise ist \(t<p+q\) genau dann unter einer
Summe \(a+b\) mit \(a<p,b<q\) gelegen; die nichttriviale Richtung folgt aus
der rationalen Dichtheit. Für Produkt und Kehrwert wird zusätzlich nach dem
Vorzeichen von \(p,q\) unterschieden; die rationalen Körpergesetze aus
\FormulaRefAuto{RationalOrderedFieldAxioms}[def] liefern jeweils dieselbe
Randbedingung des Hauptschnittes.

Insbesondere dürfen die rationalen Zahlen von nun an über \(\RealEmb\) mit
ihren Hauptschnitten identifiziert werden. Diese Identifikation ändert weder
Gleichungen noch Ungleichungen.

\chapter{Betrag, Abstand und Intervalle}

\section{Betrag und Abstand}

\FormulaDefDeltaK[Absolutbetrag einer reellen Zahl]{%
  A\in\mathbb R\vdash
  \RealAbs A\coloneqq
  \begin{cases}
    A,&\RealZero\RealLe A,\\
    \RealNeg A,&A\RealLt\RealZero.
  \end{cases}%
}{RealAbsoluteValue}{%
  \DeltaRow{Mengen}{A}%
  \DeltaRow{Relationsoperatoren}{\RealLe\dsep\RealLt}%
  \DeltaRow{Funktionsoperatoren}{\RealNeg{\,\cdot\,}}%
  \DeltaRow{\textbf{Neues Symbol}}{}%
  \DeltaRow{Reeller Betrag}{\RealAbs{\,\cdot\,}}%
}

\FormulaDefDeltaK[Abstand zweier reeller Zahlen]{%
  A,B\in\mathbb R\vdash
  \RealDist{A}{B}\coloneqq\RealAbs{A\RealAdd\RealNeg{B}}%
}{RealDistance}{%
  \DeltaRow{Mengen}{A\dsep B}%
  \DeltaRow{Funktionsoperatoren}{\RealAdd\dsep\RealNeg{\,\cdot\,}
    \dsep\RealAbs{\,\cdot\,}\dsep\RealDist{\,\cdot\,}{\,\cdot\,}}%
  \DeltaRow{\textbf{Neues Symbol}}{}%
  \DeltaRow{Reeller Abstand}{\RealDist{\,\cdot\,}{\,\cdot\,}}%
}

\FormulaThmDeltaK[Grundgesetze des reellen Betrags]{%
  \begin{aligned}[t]
  A,B\in\mathbb R\vdash{}&
    \RealZero\RealLe\RealAbs A\dsep
    \bigl(\RealAbs A=\RealZero\leftrightarrow A=\RealZero\bigr)\dsep{}\\[-2pt]
  &\RealAbs{\RealNeg A}=\RealAbs A\dsep
    \RealAbs{A\RealMul B}=\RealAbs A\RealMul\RealAbs B\dsep{}\\[-2pt]
  &\RealAbs{A\RealAdd B}\RealLe\RealAbs A\RealAdd\RealAbs B.
  \end{aligned}%
}{RealAbsoluteValueLaws}{%
  \DeltaRow{Mengen}{A\dsep B}%
  \DeltaRow{Relationsoperatoren}{\RealLe}%
  \DeltaRow{Funktionsoperatoren}{\RealAdd\dsep\RealMul
    \dsep\RealNeg{\,\cdot\,}\dsep\RealAbs{\,\cdot\,}}%
}

\FormulaThmDeltaK[Grundgesetze des reellen Abstands]{%
  \begin{aligned}[t]
  A,B,C\in\mathbb R\vdash{}&
    \RealZero\RealLe\RealDist{A}{B}\dsep
    \bigl(\RealDist{A}{B}=\RealZero\leftrightarrow A=B\bigr)\dsep{}\\[-2pt]
  &\RealDist{A}{B}=\RealDist{B}{A}\dsep
    \RealDist{A}{C}\RealLe
      \RealDist{A}{B}\RealAdd\RealDist{B}{C}.
  \end{aligned}%
}{RealDistanceLaws}{%
  \DeltaRow{Mengen}{A\dsep B\dsep C}%
  \DeltaRow{Relationsoperatoren}{\RealLe}%
  \DeltaRow{Funktionsoperatoren}{\RealAdd\dsep\RealDist{\,\cdot\,}{\,\cdot\,}}%
}

\section{Intervalle und Umgebungen}

\FormulaDefDeltaK[Beschränkte reelle Intervalle]{%
  \begin{aligned}[t]
  A,B\in\mathbb R\vdash\quad
  [A,B]_{\mathbb R}&\coloneqq
    \{x\in\mathbb R\mid A\RealLe x\land x\RealLe B\},\\
  (A,B)_{\mathbb R}&\coloneqq
    \{x\in\mathbb R\mid A\RealLt x\land x\RealLt B\},\\
  [A,B)_{\mathbb R}&\coloneqq
    \{x\in\mathbb R\mid A\RealLe x\land x\RealLt B\},\\
  (A,B]_{\mathbb R}&\coloneqq
    \{x\in\mathbb R\mid A\RealLt x\land x\RealLe B\}.
  \end{aligned}%
}{RealBoundedIntervals}{%
  \DeltaRow{Mengen}{A\dsep B\dsep x}%
  \DeltaRow{Relationsoperatoren}{\RealLe\dsep\RealLt}%
  \DeltaRow{\textbf{Neue Symbole}}{}%
  \DeltaRow{Beschränkte Intervalle}
    {[A,B]_{\mathbb R}\dsep(A,B)_{\mathbb R}
      \dsep[A,B)_{\mathbb R}\dsep(A,B]_{\mathbb R}}%
}

\FormulaDefDeltaK[Unbeschränkte reelle Intervalle]{%
  \begin{aligned}[t]
  A,B\in\mathbb R\vdash\quad
  (-\infty,B)_{\mathbb R}&\coloneqq
    \{x\in\mathbb R\mid x\RealLt B\},\\
  (-\infty,B]_{\mathbb R}&\coloneqq
    \{x\in\mathbb R\mid x\RealLe B\},\\
  (A,\infty)_{\mathbb R}&\coloneqq
    \{x\in\mathbb R\mid A\RealLt x\},\\
  [A,\infty)_{\mathbb R}&\coloneqq
    \{x\in\mathbb R\mid A\RealLe x\}.
  \end{aligned}%
}{RealUnboundedIntervals}{%
  \DeltaRow{Mengen}{A\dsep B\dsep x}%
  \DeltaRow{Relationsoperatoren}{\RealLe\dsep\RealLt}%
  \DeltaRow{\textbf{Neue Symbole}}{}%
  \DeltaRow{Unbeschränkte Intervalle}
    {(-\infty,B)_{\mathbb R}\dsep(-\infty,B]_{\mathbb R}}%
  \DeltaRow{Unbeschränkte Intervalle}
    {(A,\infty)_{\mathbb R}\dsep[A,\infty)_{\mathbb R}}%
}

\FormulaDefDeltaK[Offene Epsilon-Umgebung]{%
  \begin{aligned}[t]
  &x,\varepsilon\in\mathbb R\dsep\RealZero\RealLt\varepsilon\\[-2pt]
  &\quad\vdash
  B_{\varepsilon}(x)\coloneqq
    \{y\in\mathbb R\mid\RealDist{y}{x}\RealLt\varepsilon\}.
  \end{aligned}%
}{RealEpsilonNeighborhood}{%
  \DeltaRow{Mengen}{x\dsep y\dsep\varepsilon}%
  \DeltaRow{Relationsoperatoren}{\RealLt}%
  \DeltaRow{Funktionsoperatoren}{\RealDist{\,\cdot\,}{\,\cdot\,}}%
  \DeltaRow{\textbf{Neues Symbol}}{}%
  \DeltaRow{Offene Umgebung}{B_{\varepsilon}(x)}%
}

\FormulaThmDeltaK[Intervallform der Epsilon-Umgebung]{%
  x,\varepsilon\in\mathbb R\dsep\RealZero\RealLt\varepsilon
  \vdash
  B_{\varepsilon}(x)=
  \bigl(x\RealAdd\RealNeg\varepsilon,
        x\RealAdd\varepsilon\bigr)_{\mathbb R}%
}{RealEpsilonNeighborhoodInterval}{%
  \DeltaRow{Mengen}{x\dsep\varepsilon}%
  \DeltaRow{Relationsoperatoren}{\RealLt}%
  \DeltaRow{Funktionsoperatoren}{\RealAdd\dsep\RealNeg{\,\cdot\,}
    \dsep\RealDist{\,\cdot\,}{\,\cdot\,}}%
}

\chapter{Archimedische Grundbegriffe}

Die kanonische reelle Einbettung einer natürlichen Zahl ist die Komposition
\[
  n\longmapsto
  \RealEmb\bigl(\RatEmbed(\IntEmbed(n))\bigr),
\]
diejenige einer ganzen Zahl ist
\(z\mapsto\RealEmb(\RatEmbed(z))\). In den folgenden Formeln werden diese
Kompositionen ausgeschrieben; damit bleiben alle Zwischentypen aus den
Bänden~13 und~14 sichtbar.

\FormulaDefDeltaK[Archimedische Einbettung]{%
  \begin{aligned}[t]
  &\eta\colon\mathbb N\to K\dsep\TotOrd{K,\preccurlyeq}\\[-2pt]
  &\quad\vdash
  \operatorname{Arch}(K,\preccurlyeq,\eta)
  \coloneqq
  \forall x\in K\,\exists n\in\mathbb N\,
    \bigl(x\preccurlyeq\eta(n)\land x\neq\eta(n)\bigr).
  \end{aligned}%
}{ArchimedeanOrderedEmbedding}{%
  \DeltaRow{Mengen}{K\dsep x\dsep n}%
  \DeltaRow{Funktionen}{\eta}%
  \DeltaRow{Relationsoperatoren}{\preccurlyeq}%
  \DeltaRow{\textbf{Neues Symbol}}{}%
  \DeltaRow{Archimedisches Prädikat}
    {\operatorname{Arch}(K,\preccurlyeq,\eta)}%
}

\FormulaThmDeltaK[Archimedische Eigenschaft der reellen Zahlen]{%
  \forall x\in\mathbb R\,\exists n\in\mathbb N\,
    \bigl(x\RealLt
      \RealEmb(\RatEmbed(\IntEmbed(n)))\bigr)%
}{RealArchimedeanProperty}{%
  \DeltaRow{Mengen}{x\dsep n}%
  \DeltaRow{Funktionen}{\IntEmbed\dsep\RatEmbed\dsep\RealEmb}%
  \DeltaRow{Relationsoperatoren}{\RealLt}%
}

\emph{Beweisskizze.}
Sei \(x=A\in\mathbb R\). Wegen \(A\neq\mathbb Q\) wähle
\(q\in\mathbb Q\setminus A\). Nach
\FormulaRefAuto{RationalArchimedean}[thm] gibt es \(n\in\mathbb N\) mit
\(q<\RatEmbed(\IntEmbed(n))\). Jedes \(a\in A\) erfüllt \(a\leq q\), denn
aus \(q<a\) folgte durch Abwärtsabschluss \(q\in A\). Also liegt \(A\) im
entsprechenden Hauptschnitt; die Inklusion ist echt, weil \(q\) dem
Hauptschnitt, aber nicht \(A\) angehört.

\FormulaThmDeltaK[Ganzzahliger Abschnitt einer reellen Zahl]{%
  \forall x\in\mathbb R\,\exists z\in\mathbb Z\,
  \bigl(\RealEmb(\RatEmbed(z))\RealLe x\land
        x\RealLt\RealEmb(\RatEmbed(z+1_{\mathbb Z}))\bigr)%
}{RealIntegerPartProperty}{%
  \DeltaRow{Mengen}{x\dsep z}%
  \DeltaRow{Funktionen}{\RatEmbed\dsep\RealEmb}%
  \DeltaRow{Relationsoperatoren}{\RealLe\dsep\RealLt}%
}

\emph{Beweisskizze.}
Nach der archimedischen Eigenschaft sind die ganzen Zahlen oberhalb und,
nach Anwendung auf \(-x\), unterhalb von \(x\) nichtleer. Wähle eine ganze
untere Schranke \(w\) und betrachte die natürlichen \(n\), für die
\(x\RealLt\RealEmb(\RatEmbed(w+\IntEmbed(n)))\) gilt. Diese Menge ist
nichtleer und besitzt nach dem Wohlordnungsprinzip aus Band~12 ein kleinstes
Element; es ist wegen der unteren Schranke ein Nachfolger \(\Succ(k)\).
Mit \(z=w+\IntEmbed(k)\) liefern die Minimalität und die Ordnungstreue der
Einbettungen genau
\(\RealEmb(\RatEmbed(z))\RealLe x
\RealLt\RealEmb(\RatEmbed(z+1_{\mathbb Z}))\).

\FormulaThmDeltaK[Dichtheit der rationalen Zahlen in den reellen Zahlen]{%
  x,y\in\mathbb R\dsep x\RealLt y
  \vdash
  \exists q\in\mathbb Q\,
    \bigl(x\RealLt\RealEmb(q)\land\RealEmb(q)\RealLt y\bigr)%
}{RationalDenseInReal}{%
  \DeltaRow{Mengen}{x\dsep y\dsep q}%
  \DeltaRow{Funktionen}{\RealEmb}%
  \DeltaRow{Relationsoperatoren}{\RealLt}%
}

\emph{Beweisskizze.}
Schreibe \(x=A\subsetneq B=y\) und wähle \(b\in B\setminus A\). Da \(B\)
kein größtes Element besitzt, gibt es \(q\in B\) mit \(b<q\). Dann gilt
\(A\subsetneq\RealHat q\): Jedes Element von \(A\) liegt höchstens bei \(b\)
und damit unter \(q\), während \(b\notin A\) im Hauptschnitt liegt. Ferner
ist \(\RealHat q\subsetneq B\), weil \(B\) abwärts abgeschlossen ist und
\(q\in B\setminus\RealHat q\). Mit
\FormulaRefAuto{RationalToRealEmbedding}[def] folgt die Behauptung.

\FormulaThmDeltaK[Beliebig kleine positive rationale Hauptschnitte]{%
  \begin{aligned}[t]
  &\varepsilon\in\mathbb R\dsep\RealZero\RealLt\varepsilon
    \vdash \exists n\in\mathbb N\,\bigl(0<n\land{}\\[-2pt]
  &\qquad \RealZero\RealLt
      \RealEmb\bigl((\RatEmbed(\IntEmbed(n)))^{-1}\bigr)
    \land{}\\[-2pt]
  &\qquad \RealEmb\bigl((\RatEmbed(\IntEmbed(n)))^{-1}\bigr)
        \RealLt\varepsilon\bigr).
  \end{aligned}%
}{RealSmallPositiveRational}{%
  \DeltaRow{Mengen}{\varepsilon\dsep n}%
  \DeltaRow{Funktionen}{\IntEmbed\dsep\RatEmbed\dsep\RealEmb}%
  \DeltaRow{Relationsoperatoren}{\RealLt}%
}

\emph{Beweisskizze.}
Wende \FormulaRefAuto{RealArchimedeanProperty}[thm] auf
\(\RealInv\varepsilon\) an und ersetze den gefundenen natürlichen Index
gegebenenfalls durch seinen Nachfolger, so dass er positiv ist.
Ordnungskompatibilität und
die Kehrwertgesetze ergeben nach Multiplikation mit den positiven Größen
\(\varepsilon\) und
\(\RealEmb(\RatEmbed(\IntEmbed(n)))\) die Ungleichung
\[
  \RealZero\RealLt
  \RealEmb\bigl((\RatEmbed(\IntEmbed(n)))^{-1}\bigr)
  \RealLt\varepsilon .
\]

Die letzte Aussage ist die Form der archimedischen Eigenschaft, die in
späteren Epsilon-Argumenten am häufigsten verwendet wird. Sie setzt nur die
positive reelle Zahl \(\varepsilon\) voraus und liefert eine rationale
Vergleichsgröße innerhalb derselben Zielmenge \(\mathbb R\).

\chapter{Axiomatische Zusammenfassung}

\section{Vollständig angeordnete Körper}

Die folgende Definition fasst genau die Eigenschaften zusammen, welche die
Schnittkonstruktion bereitstellt. Die Vollständigkeitszeile quantifiziert über
Teilmengen des Trägers und ist daher in dieser Form zweitstufig zu lesen.

\FormulaDefDeltaK[Vollständig angeordneter Körper]{%
  \CompleteOrderedField{\begin{gathered}
    K,\boxplus,\boxtimes,\boxminus,\\[-2pt]
    \operatorname{inv},0_K,1_K,\preccurlyeq
  \end{gathered}}%
}{CompleteOrderedFieldAxioms}{%
  \DeltaRow{Mengen}{K\dsep x\dsep y\dsep z}%
  \DeltaRow{Mengen}{S\dsep s\dsep u}%
  \DeltaRow{Relationsoperatoren}{\preccurlyeq}%
  \DeltaRow{Funktionsoperatoren}{\boxplus\dsep\boxtimes}%
  \DeltaRow{Funktionsoperatoren}{\boxminus\dsep\operatorname{inv}}%
  \DeltaRow{\textbf{Typen}}{}%
  \DeltaRow{Operationen}
    {\boxplus,\boxtimes\colon K\times K\to K}%
  \DeltaRow{Negation}{\boxminus\colon K\to K}%
  \DeltaRow{Kehrwert}
    {\operatorname{inv}\colon K\setminus\{0_K\}\to K}%
  \DeltaRow{Konstanten}{0_K,1_K\in K\dsep0_K\neq1_K}%
  \DeltaRow{\textbf{Körpergesetze}}{}%
  \DeltaRow{Ass. Addition}
    {x,y,z\in K\vdash
      (x\boxplus y)\boxplus z=x\boxplus(y\boxplus z)}%
  \DeltaRow{Komm. Addition}
    {x,y\in K\vdash x\boxplus y=y\boxplus x}%
  \DeltaRow{Null und Inverses}
    {x\in K\vdash x\boxplus0_K=x\dsep
      x\boxplus(\boxminus x)=0_K}%
  \DeltaRow{Ass. Produkt}
    {x,y,z\in K\vdash
      (x\boxtimes y)\boxtimes z=x\boxtimes(y\boxtimes z)}%
  \DeltaRow{Komm. Produkt}
    {x,y\in K\vdash x\boxtimes y=y\boxtimes x}%
  \DeltaRow{Eins}
    {x\in K\vdash x\boxtimes1_K=x}%
  \DeltaRow{Produktinverses}
    {x\in K\setminus\{0_K\}\vdash
      x\boxtimes\operatorname{inv}(x)=1_K}%
  \DeltaRow{Distributivität}
    {x,y,z\in K\vdash x\boxtimes(y\boxplus z)=
      (x\boxtimes y)\boxplus(x\boxtimes z)}%
  \DeltaRow{\textbf{Ordnung}}{}%
  \DeltaRow{Totalität}{\TotOrd{K,\preccurlyeq}}%
  \DeltaRow{Translation}
    {x,y,z\in K\dsep x\preccurlyeq y\vdash
      x\boxplus z\preccurlyeq y\boxplus z}%
  \DeltaRow{Positive Produkte}
    {x,y\in K\dsep0_K\preccurlyeq x\dsep0_K\preccurlyeq y
      \vdash0_K\preccurlyeq x\boxtimes y}%
  \DeltaRow{Vollständigkeit}
    {\begin{aligned}[t]
      &\varnothing\neq S\subseteq K\dsep
        \UB_{K,\preccurlyeq}(S)\neq\varnothing\vdash{}\\[-2pt]
      &\exists s\in\UB_{K,\preccurlyeq}(S)\,
        \forall u\in\UB_{K,\preccurlyeq}(S)\,(s\preccurlyeq u)
     \end{aligned}}%
  \DeltaRow{\textbf{Struktur}}{}%
  \DeltaRow{Prädikat}{\mathsf{COF}}%
}

\FormulaThmDeltaK[Das Schnittmodell ist ein vollständig angeordneter Körper]{%
  \CompleteOrderedField{\mathbb R,\RealAdd,\RealMul,
    \RealNeg{\,\cdot\,},\RealInv{\,\cdot\,},
    \RealZero,\RealOne,\RealLe}%
}{RealCutsCompleteOrderedField}{%
  \DeltaRow{Mengen}{\mathbb R}%
  \DeltaRow{Relationsoperatoren}{\RealLe}%
  \DeltaRow{Funktionsoperatoren}{\RealAdd\dsep\RealMul
    \dsep\RealNeg{\,\cdot\,}\dsep\RealInv{\,\cdot\,}}%
}

\emph{Beweisskizze.}
Die Typisierung liefern
\FormulaRefAuto{RealCutAdditiveClosure}[thm] und
\FormulaRefAuto{RealCutMultiplicativeClosure}[thm]. Null und Eins sind
Hauptschnitte; sie sind verschieden, weil
\(0_{\mathbb Q}\neq1_{\mathbb Q}\) und die rationale Einbettung injektiv ist.
Die Körper- und Ordnungszeilen sind
\FormulaRefAuto{RealCutAdditiveLaws}[thm],
\FormulaRefAuto{RealCutMultiplicativeLaws}[thm],
\FormulaRefAuto{RealCutTotalOrder}[thm] und
\FormulaRefAuto{RealCutOrderArithmeticCompatibility}[thm]. Die
Vollständigkeitszeile ist genau die Schrankenmengenform von
\FormulaRefAuto{RealLeastUpperBoundProperty}[thm].

\FormulaThmDeltaK[Kanonische natürliche Abbildung eines angeordneten Körpers]{%
  \begin{aligned}[t]
  &\CompleteOrderedField{K,\boxplus,\boxtimes,
    \boxminus,\operatorname{inv},0_K,1_K,\preccurlyeq}\\[-2pt]
  &\quad\vdash
    \exists!\eta\colon\mathbb N\to K\,
    \left(
      \eta(0)=0_K\land
      \forall n\in\mathbb N\,
        \bigl(\eta(\Succ(n))=\eta(n)\boxplus1_K\bigr)
    \right).
  \end{aligned}%
}{CompleteOrderedFieldNaturalEmbedding}{%
  \DeltaRow{Mengen}{K\dsep n}%
  \DeltaRow{Funktionen}{\eta}%
  \DeltaRow{Funktionsoperatoren}{\boxplus}%
}

\emph{Beweisskizze.}
Wegen \(0_K\in K\) und \(x\mapsto x\boxplus1_K\colon K\to K\) ist dies
unmittelbar der Dedekindsche Rekursionssatz aus Band~12. Wir schreiben
\(\eta_K\) für die eindeutige Funktion. Nur diese durch die beiden
Rekursionsgleichungen charakterisierte kanonische Abbildung ist im folgenden
Satz gemeint.

\FormulaThmDeltaK[Archimedizität vollständig angeordneter Körper]{%
  \begin{aligned}[t]
  &\CompleteOrderedField{\begin{gathered}
      K,\boxplus,\boxtimes,\boxminus,\\[-2pt]
      \operatorname{inv},0_K,1_K,\preccurlyeq
    \end{gathered}}\dsep{}\\[-2pt]
  &\eta_K\colon\mathbb N\to K\dsep
    \eta_K(0)=0_K\dsep{}\\[-2pt]
  &\forall n\in\mathbb N\,
      \bigl(\eta_K(\Succ(n))=\eta_K(n)\boxplus1_K\bigr)\\[-2pt]
  &\quad\vdash
    \operatorname{Arch}(K,\preccurlyeq,\eta_K).
  \end{aligned}%
}{CompleteOrderedFieldArchimedean}{%
  \DeltaRow{Mengen}{K\dsep n}%
  \DeltaRow{Funktionen}{\eta_K}%
  \DeltaRow{Relationsoperatoren}{\preccurlyeq}%
}

\emph{Beweisskizze.}
Wäre \(\eta_K[\mathbb N]\) nach oben beschränkt, so besäße diese
nichtleere Menge ein Supremum \(s\). Aus den angeordneten Körperaxiomen folgt
\(0_K\preccurlyeq1_K\) und \(0_K\neq1_K\). Daher ist
\(s\boxplus(\boxminus1_K)\) keine obere Schranke, und es gibt \(n\) mit
\[
  s\boxplus(\boxminus1_K)\preccurlyeq\eta_K(n)
  \quad\text{und}\quad
  s\boxplus(\boxminus1_K)\neq\eta_K(n).
\]
Translation um \(1_K\) und die Rekursionsgleichung liefern
\(s\preccurlyeq\eta_K(\Succ(n))\) bei Verschiedenheit, im Widerspruch zur
Oberen-Schranken-Eigenschaft von \(s\). Also ist das Bild der kanonischen
Abbildung unbeschränkt.

\FormulaThmDeltaK[Eindeutigkeit des vollständig angeordneten reellen Körpers]{%
  \begin{aligned}[t]
  &\CompleteOrderedField{K,\boxplus,\boxtimes,
    \boxminus,\operatorname{inv},0_K,1_K,\preccurlyeq}\\[-2pt]
  &\quad\vdash\exists!\Phi\colon\mathbb R\bij K\,
    \left(
      \begin{aligned}[t]
      &\forall x,y\in\mathbb R\,
        \bigl(\Phi(x\RealAdd y)=\Phi(x)\boxplus\Phi(y)\bigr)\land{}\\[-2pt]
      &\forall x,y\in\mathbb R\,
        \bigl(\Phi(x\RealMul y)=\Phi(x)\boxtimes\Phi(y)\bigr)\land{}\\[-2pt]
      &\forall x,y\in\mathbb R\,
        \bigl(x\RealLe y\leftrightarrow\Phi(x)\preccurlyeq\Phi(y)\bigr)\land{}\\[-2pt]
      &\Phi(\RealZero)=0_K\land\Phi(\RealOne)=1_K
      \end{aligned}
    \right).
  \end{aligned}%
}{CompleteOrderedFieldUniqueness}{%
  \DeltaRow{Mengen}{K\dsep x\dsep y}%
  \DeltaRow{Funktionen}{\Phi}%
  \DeltaRow{Relationsoperatoren}{\RealLe\dsep\preccurlyeq}%
  \DeltaRow{Funktionsoperatoren}{\RealAdd\dsep\RealMul
    \dsep\boxplus\dsep\boxtimes}%
}

\emph{Beweisskizze.}
Die kanonische natürliche Abbildung erzeugt mit Negation und Kehrwert die
eindeutige ordnungserhaltende Einbettung
\(\eta_{\mathbb Q}\colon\mathbb Q\to K\) des rationalen Primkörpers. Für
einen Dedekind-Schnitt \(A\) ist \(\eta_{\mathbb Q}[A]\) nichtleer und durch
jedes \(\eta_{\mathbb Q}(r)\) mit \(r\notin A\) nach oben beschränkt. Setze
deshalb
\[
  \Phi(A)\coloneqq
  \sup_{K,\preccurlyeq}\bigl(\eta_{\mathbb Q}[A]\bigr).
\]
Die Vollständigkeitszeile des Schemas macht diese Vorschrift wohldefiniert.
Aus \(A\subsetneq B\) wählt man \(q\in B\setminus A\) und anschließend wegen
des fehlenden größten Elements ein \(r\in B\) mit \(q<r\). Damit liegt das
Supremum von \(B\) strikt über dem von \(A\); \(\Phi\) ist also
ordnungserhaltend und injektiv.

Für \(y\in K\) bildet
\[
  A_y\coloneqq
  \{q\in\mathbb Q\mid
    \eta_{\mathbb Q}(q)\preccurlyeq y
    \land\eta_{\mathbb Q}(q)\neq y\}
\]
einen Dedekind-Schnitt. Nichtleerheit, Echtheit und das Fehlen eines größten
Elements folgen aus
\FormulaRefAuto{CompleteOrderedFieldArchimedean}[thm] und der daraus durch den
üblichen Nennervergleich gewonnenen Dichtheit des rationalen Primkörpers in
\(K\). Die Supremumseigenschaft gibt \(\Phi(A_y)=y\), also Surjektivität.
Die rationalen Rechengesetze und die Verträglichkeit von Suprema mit
Translation und positiver Skalierung liefern die beiden Operationsgleichungen.
Schließlich muss jeder ordnungs- und körpererhaltende Isomorphismus
\(\eta_{\mathbb Q}\) festhalten und Suprema bewahren; die angezeigte
Konstruktion ist daher die einzig mögliche.

Die Eindeutigkeit ist als Eindeutigkeit des ordnungs- und
körperstrukturerhaltenden Isomorphismus zu verstehen. Damit kann in späteren
Bänden entweder mit dem konkreten Schnittmodell oder ausschließlich mit den
registrierten Axiomen eines vollständig angeordneten Körpers gearbeitet
werden. Für Grenzwertaussagen stehen insbesondere die reelle Ordnung, der
Abstand und positive Epsilon-Umgebungen bereits zur Verfügung.

\end{document}
